\documentclass[12pt,a4paper]{report}

% ==========================================
% PACKAGES ET CONFIGURATION
% ==========================================
\usepackage[utf8]{inputenc}
\usepackage[T1]{fontenc}
\usepackage[french]{babel}
\usepackage{graphicx}
\usepackage[dvipsnames, table]{xcolor}
\usepackage{float}
\usepackage{tikz}
\usepackage{tcolorbox}
\usepackage{xparse}
\usepackage{amssymb}

% ==========================================
% CONFIGURATION DES MARGES (ISET Siliana)
% ==========================================
\usepackage{geometry}
\geometry{
  a4paper,
  top=2cm,
  bottom=2cm,
  left=2.5cm,        % 2cm + 0.5cm reliure
  right=2cm,
  headsep=0.75cm,
  headheight=1.25cm,
  footskip=1.25cm
}

% ==========================================
% POLICE DE CARACTÈRES
% ==========================================
\usepackage{mathptmx}  % Times New Roman
\usepackage{setspace}
\onehalfspacing        % Interligne 1.5 ligne

% ==========================================
% EN-TÊTES ET PIEDS DE PAGE
% ==========================================
\usepackage{fancyhdr}
\pagestyle{fancy}
\fancyhf{}

% En-tête personnalisé
\fancyhead[C]{%
  \fontsize{10}{12}\selectfont\itshape%
  Farm AI : Système Intelligent de Gestion Agricole%
}
\renewcommand{\headrulewidth}{0.5pt}

% Pied de page personnalisé
\fancyfoot[L]{\fontsize{10}{12}\selectfont\itshape Iset Siliana}
\fancyfoot[C]{\fontsize{10}{12}\selectfont\itshape \thepage}
\fancyfoot[R]{\fontsize{10}{12}\selectfont\itshape Luniweb / Farm AI}
\renewcommand{\footrulewidth}{0.5pt}

% ==========================================
% STYLE DES TITRES
% ==========================================
\usepackage{titlesec}

% Chapitre : 20pt, Gras
\titleformat{\chapter}[block]
  {\normalfont\fontsize{20}{24}\selectfont\bfseries}
  {\chaptertitlename\ \thechapter\ :}
  {0.5em}
  {}
\titlespacing*{\chapter}{0pt}{12pt}{35pt}

% Section : 16pt, Gras
\titleformat{\section}
  {\normalfont\fontsize{16}{20}\selectfont\bfseries\color{black}}
  {\thesection.}
  {1em}
  {}
\titlespacing*{\section}{20pt}{20pt}{10pt}

% Subsection : 14pt, Gras
\titleformat{\subsection}
  {\normalfont\fontsize{14}{18}\selectfont\bfseries}
  {\thesubsection.}
  {1em}
  {}
\titlespacing*{\subsection}{10pt}{10pt}{5pt}

% Subsubsection : 12pt, Gras + Italique
\renewcommand{\thesubsubsection}{\alph{subsubsection}.}
\titleformat{\subsubsection}
  {\normalfont\fontsize{12}{16}\selectfont\bfseries\itshape}
  {\thesubsubsection}
  {1em}
  {}
\titlespacing*{\subsubsection}{5pt}{5pt}{0pt}

% ==========================================
% TABLEAUX AVANCÉS
% ==========================================
\usepackage{booktabs}
\usepackage{array}
\usepackage{tabularx}
\usepackage{colortbl}

% Commande personnalisée pour tableaux
\newcommand{\tableheader}[1]{
  \cellcolor[gray]{0.8}%
  \bfseries#1%
}

% ==========================================
% FIGURES ET LÉGENDES
% ==========================================
\usepackage{caption}
\captionsetup{
  font={it,small},
  justification=centering,
  labelsep=colon,
  format=plain
}

% Numérotation des figures/tableaux par chapitre
\numberwithin{figure}{chapter}
\numberwithin{table}{chapter}

% ==========================================
% LISTES ET ÉNUMÉRATIONS
% ==========================================
\usepackage{enumitem}
\setlist[itemize]{
  leftmargin=1.5cm,
  itemsep=0.5mm,
  parsep=0mm,
  topsep=2mm
}
\setlist[enumerate]{
  leftmargin=1.5cm,
  itemsep=0.5mm,
  parsep=0mm,
  topsep=2mm
}

% ==========================================
% LIENS HYPERTEXTE
% ==========================================
\usepackage{hyperref}
\hypersetup{
  colorlinks=true,
  linkcolor=black,
  filecolor=black,
  urlcolor=blue,
  citecolor=black,
  pdftitle={Farm AI - Rapport PFE},
  pdfauthor={Mohamed Amine ABBASSI},
  pdfsubject={Système Intelligent de Gestion Agricole}
}

% ==========================================
% COMMANDES PERSONNALISÉES
% ==========================================

% Box d'information
\newtcolorbox{infobox}{
  colback=blue!5!white,
  colframe=blue!50!black,
  rounded corners,
  fonttitle=\bfseries,
  title=ℹ Information
}

% Box d'attention
\newtcolorbox{warningbox}{
  colback=orange!5!white,
  colframe=orange!75!black,
  rounded corners,
  fonttitle=\bfseries,
  title=⚠ Attention
}

% Chaque chapitre commence sur une nouvelle page
\usepackage{etoolbox}
\pretocmd{\chapter}{\clearpage}{}{}

% Reférencement automatique des figures
\graphicspath{{./image/}{./images/}{./figures/}}

% ==========================================
% DÉBUT DU DOCUMENT
% ==========================================
\begin{document}

% Page de garde sans en-tête ni pied de page
\thispagestyle{empty}

% ===========================================================
% PAGE DE GARDE (Modèle ISET Siliana)
% ===========================================================
\begin{titlepage}
  \begin{center}

    % En-tête institutionnel
    {\large \textbf{REPUBLIQUE TUNISIENNE}} \\[2pt]
    {\normalsize \textbf{MINISTERE DE L'ENSEIGNEMENT SUPERIEUR}} \\
    {\normalsize \textbf{ET DE LA RECHERCHE SCIENTIFIQUE}} \\[2pt]
    {\normalsize \textbf{DIRECTION GENERALE DES ETUDES TECHNOLOGIQUES}} \\[0.8cm]

    % Logo et établissement
    {\Large \textbf{Institut Supérieur des Études Technologiques de Siliana}} \\[4pt]
    {\large \textbf{Département de Génie Électrique}}

    \vspace{1.5cm}
    \rule{\linewidth}{0.8pt}\\[0.4cm]
    {\Large \textbf{Projet de fin d'études}}\\[0.4cm]
    \rule{\linewidth}{0.8pt}

    \vspace{1cm}

    % Titre du projet
    {\LARGE \textbf{Farm AI}}\\[0.3cm]
    {\LARGE \textbf{Système intelligent de gestion agricole}}

    \vspace{2cm}

    % Bloc étudiant / encadrants
    \begin{minipage}{0.45\textwidth}
      \begin{flushleft}
        \emph{Élaboré par :}\\[6pt]
        \textbf{Mohamed Amine ABBASSI}
      \end{flushleft}
    \end{minipage}
    \hfill
    \begin{minipage}{0.45\textwidth}
      \begin{flushright}
        \emph{Encadré par :}\\[4pt]
        \textbf{Mme. Abir KHALDI} (ISET)\\
        \textbf{Mr. Rochdi ABID} (Luniweb)\\[8pt]
        \emph{Entreprise d'accueil :}\\[4pt]
        \textbf{Luniweb} (Agence Web)
      \end{flushright}
    \end{minipage}

    \vfill

    {\large \textbf{Année universitaire : 2025 -- 2026}}\\[4pt]
    {\normalsize Promotion : 2025 -- Semestre 2}

  \end{center}
\end{titlepage}

% ===========================================================
% DÉDICACE
% ===========================================================
\chapter*{Dédicace}
\addcontentsline{toc}{chapter}{Dédicace}

\begin{center}
  \textit{\large
    Je dédie ce travail à ma famille et à tous ceux qui m'ont soutenu\\[0.5cm]
    tout au long de mon parcours académique et professionnel.
  }
\end{center}

\vspace{3cm}

% ===========================================================
% REMERCIEMENTS
% ===========================================================
\chapter*{Remerciements}
\addcontentsline{toc}{chapter}{Remerciements}

\begin{itshape}
  Au seuil de ce travail, nous souhaitons exprimer notre profonde gratitude à nos encadreurs,
  Madame Abir KHALDI et Monsieur Rochdi ABID, pour la qualité de leur accompagnement, leur
  disponibilité constante ainsi que la pertinence de leurs conseils tout au long de la réalisation
  de ce projet.

  Nous adressons également nos remerciements les plus respectueux aux membres du jury qui nous ont
  fait l'honneur d'accepter d'évaluer ce travail.

  Enfin, nous tenons à remercier l'équipe de Luniweb pour son accueil, sa confiance et son soutien
  tout au long de cette période de stage.
\end{itshape}

% ===========================================================
% TABLE DES MATIÈRES
% ===========================================================
\newpage
\tableofcontents

% ===========================================================
% LISTES DES FIGURES ET TABLEAUX
% ===========================================================
\newpage
\listoffigures

\newpage
\listoftables

% ===========================================================
% INTRODUCTION GÉNÉRALE
% ===========================================================
\chapter*{Introduction générale}
\addcontentsline{toc}{chapter}{Introduction générale}

\section*{Contexte général}
L'intégration de l'Intelligence Artificielle dans le domaine agricole représente aujourd'hui une
évolution majeure des pratiques de gestion. Dans un contexte où la traçabilité du personnel et la
sécurité des sites sont des enjeux critiques, les systèmes manuels montrent rapidement leurs limites.

\begin{infobox}
  L'agriculture connectée s'impose comme une nécessité dans la transformation numérique
  des exploitations modernes.
\end{infobox}

\section*{Objectifs du travail}
Ce projet a pour objectif de concevoir et de développer la plateforme \textbf{\textit{Farm AI}}, un
système moderne visant à :

\begin{itemize}
  \item Automatiser la gestion du personnel ;
  \item Assurer le suivi des présences en temps réel ;
  \item Renforcer la sécurité des sites agricoles ;
  \item Intégrer les technologies de vision par ordinateur et d'IA.
\end{itemize}

\section*{Organisation du rapport}
Le présent rapport est structuré en quatre chapitres :

\begin{enumerate}
  \item \textbf{Contexte et État de l'art} : présentation du contexte et étude comparative ;
  \item \textbf{Spécification des besoins et Conception} : analyse des besoins et modélisation UML ;
  \item \textbf{Environnement et Choix Techniques} : technologies et architecture ;
  \item \textbf{Réalisation et Tests} : présentation des interfaces et résultats de validation.
\end{enumerate}

% ===========================================================
% CHAPITRE 1 : CONTEXTE ET ÉTAT DE L'ART
% ===========================================================
\chapter{Contexte et État de l'art}

\section{Introduction}
Ce chapitre présente le contexte général du projet et la problématique liée à la gestion classique
des fermes, ainsi qu'une étude comparative des solutions existantes.

\section{Présentation du cadre du projet}

\subsection{Présentation de l'entreprise d'accueil}
Luniweb est une agence pluridisciplinaire spécialisée dans le numérique et le développement web,
fondée en 2015 en Tunisie. L'agence intervient dans les domaines suivants :

\begin{itemize}
  \item Développement web et mobile ;
  \item Intelligence Artificielle et Machine Learning ;
  \item Consultation en transformation numérique ;
  \item Infrastructure Cloud et DevOps.
\end{itemize}

\subsection{Contexte général du projet}
La ferme connectée est au cœur de la transition technologique actuelle. Farm AI s'inscrit dans cette
vision en combinant l'IA, la vision par ordinateur et les technologies web pour répondre aux enjeux
de gestion agricole moderne.

\section{Étude de l'existant et problématique}

\subsection{Critique de l'existant}
Les systèmes de pointage classiques et la gestion manuelle du personnel présentent de nombreuses
failles :

\begin{warningbox}
  \begin{itemize}
    \item Erreurs de saisie et oublis de pointage ;
    \item Risques de fraude et de manipulation des données ;
    \item Absence de suivi en temps réel ;
    \item Coûts administratifs et RH élevés ;
    \item Manque de traçabilité pour la sécurité.
  \end{itemize}
\end{warningbox}

\subsection{Problématique}
\textbf{Comment assurer un suivi précis, sécurisé et en temps réel des employés agricoles sans
intervention manuelle lourde ?}

\\section{Solution proposée : Farm AI}
Farm AI combine l'IA et le Web pour un pointage entièrement automatisé, offrant une solution
robuste, sécurisée et facile à déployer.

\section{Conclusion}\nCe chapitre a permis d\'identifier les lacunes des systèmes existants et de justifier le choix d\'une solution basée sur l\'intelligence artificielle. Le chapitre suivant présentera l\'état de l\'art des technologies utilisées.\n\n% ===========================================================\n% CHAPITRE 2 : ÉTAT DE L\'ART\n% ===========================================================\n\chapter{État de l\'art}
\section{Introduc tion}
Avec l’évolu tion rapide des technologies de l’intelligence ar tificielle et des systèmes intelligents, les
solutions modernes de surveillance reposent aujourd’hui sur plusieurs domaines tels que la vision par
ordinateur, le Deep Learning, le traitement des flux vidéo temps réel ainsi que les architectures web
modernes.
Dans le cadre de notre projet, plusieurs technologies et ou tils ont été u tilisés afin de concevoir une
plateforme intelligente capable d’améliorer la surveillance, le contrôle d’accès ainsi que la ges tion des
employés grâce à l’intelligence ar tificielle et à la vision par ordinateur.
Dans ce chapitre, nous présentons le domaine d’applica tion de notre projet ainsi que les principaux
concepts, architectures, technologies et ou tils utilisés pour la réalisa tion de la plateforme.
\section{Domaine de l’applica tion}
Notre projet s’inscrit dans le domaine de la surveillance intelligente et des systèmes de sécurité basés
sur l’intelligence ar tificielle. Avec l’augmenta tion des besoins en sécurité au sein des entreprises, les
systèmes tradi tionnels de surveillance devien nent insuffisants face aux nouveaux défis liés au
contrôle d’accès, à la supervision des infrastructures et à la ges tion des employés.
Les technologies modernes perme ttent aujourd’hui de développer des plateformes intelligentes
capables d’automa tiser plusieurs opéra tions de sécurité telles que la reconnaissance faciale, la
détection des personnes, l’analyse des flux vidéo ainsi que la généra tion d’alertes en temps réel.
L’objectif principal de notre projet est de proposer une solu tion intelligente perme ttant d’améliorer la
surveillance et la ges tion de la sécurité grâce aux technologies de vision par ordinateur et
d’intelligence ar tificielle.
\section{Intelligence ar tificielle et vision par ordinateur}
\subsection{Intelligence ar tificielle}
Définition
L’intelligence ar tificielle désigne l’ensemble des techniques perme ttant à une machine de simuler
certaines capacités humaines telles que l’appren tissage, l’analyse, la reconnaissance et la prise de
décision.
\textbf{Rôle dans le projet} : \\
Dans notre projet, l’intelligence ar tificielle joue un rôle essen tiel dans l’analyse des flux vidéo, la
reconnaissance faciale ainsi que la détec tion intelligente des événements de sécurité.
\textbf{Avantages} : \\
\begin{itemize}
  \item Automatisation des tâches de surveillance ;
  \item Réduction des erreurs humaines ;
  \item Analyse intelligente des données ;
  \item Amélioration de la sécurité.
\end{itemize}
\begin{figure}[H]
  \centering
  % \includegraphics[width=0.8\textwidth]{images/...}
  \caption{Schéma représentant le fonc tionnement général de l’intelligence ar tificielle.}
\end{figure}
\subsection{Vision par ordinateur}
Définition
La vision par ordinateur est un domaine de l’intelligence ar tificielle perme ttant aux machines
d’interpréter et d’analyser les images et les vidéos.
\textbf{Rôle dans le projet} : \\
Elle permet d’exploiter les flux vidéo provenant des caméras afin de détecter les personnes, analyser
les mouvements et iden tifier les visages en temps réel.
\textbf{Avantages} : \\
\begin{itemize}
  \item Analyse automa tique des vidéos ;
  \item Détection rapide des objets ;
  \item Surveillance intelligente des infrastructures.
\end{itemize}
\begin{figure}[H]
  \centering
  % \includegraphics[width=0.8\textwidth]{images/...}
  \caption{Illustra tion de la vision par ordinateur.}
\end{figure}
\subsection{Deep Learning}
Définition
Le Deep Learning est une branche du Machine Learning basée sur les réseaux de neurones profonds
permettant d’effectuer des traitements complexes sur les données.
\textbf{Rôle dans le projet} : \\
Les modèles de Deep Learning sont u tilisés pour la reconnaissance faciale, la détec tion des objets
ainsi que l’analyse intelligente des flux vidéo.
\textbf{Avantages} : \\
\begin{itemize}
  \item Haute précision ;
  \item Excellentes performances ;
  \item Apprentissage automa tique.
\end{itemize}
\begin{figure}[H]
  \centering
  % \includegraphics[width=0.8\textwidth]{images/...}
  \caption{Architecture simplifiée d’un réseau de neurones profond.}
\end{figure}
\subsection{Reconnaissance faciale}
Définition
La reconnaissance faciale est une technologie perme ttant d’iden tifier ou de vérifier l’iden tité d’une
personne à par tir de son visage.
\textbf{Rôle dans le projet} : \\
Dans notre plateforme, elle permet :
\begin{itemize}
  \item l’authentification des employés ;
  \item le contrôle d’accès intelligent ;
  \item la détection des personnes inconnues ;
  \item le suivi des entrées et sor ties.
\end{itemize}
\textbf{Avantages} : \\
\begin{itemize}
  \item Sécurité renforcée ;
  \item Authentification rapide ;
  \item Réduction des accès non autorisés.
\end{itemize}
\begin{figure}[H]
  \centering
  % \includegraphics[width=0.8\textwidth]{images/...}
  \caption{Exemple d’un système de reconnaissance faciale.}
\end{figure}
\section{Technologies de détec tion et d’analyse}
\subsection{YOLOv8}
Définition
YOLOv8 (You Only Look Once) est un modèle de Deep Learning spécialisé dans la détec tion d’objets
en temps réel.
\textbf{Rôle dans le projet} : \\
YOLOv8 permet de détecter les personnes présentes dans les flux vidéo afin d’améliorer la
surveillance intelligente des infrastructures.
\textbf{Avantages} : \\
\begin{itemize}
  \item Rapidité de traitement ;
  \item Haute précision ;
  \item Fonctionnement en temps réel.
\end{itemize}
\begin{figure}[H]
  \centering
  % \includegraphics[width=0.8\textwidth]{images/...}
  \caption{Logo YOLOv8.}
\end{figure}
\subsection{OpenCV}
Définition
OpenCV est une bibliothèque open source u tilisée dans le domaine de la vision par ordinateur et du
traitement d’images.
\textbf{Rôle dans le projet} : \\
OpenCV est u tilisé pour :
\begin{itemize}
  \item capturer les flux vidéo ;
  \item traiter les images ;
  \item prétraiter les données avant leur analyse par les modèles d’intelligence ar tificielle.
\end{itemize}
\textbf{Avantages} : \\
\begin{itemize}
  \item Compatible avec Python ;
  \item Très performante ;
  \item Large communauté.
\end{itemize}
\begin{figure}[H]
  \centering
  % \includegraphics[width=0.8\textwidth]{images/...}
  \caption{Logo OpenCV.}
\end{figure}
\subsection{InsightFace}
Définition
InsightFace est une bibliothèque spécialisée dans la reconnaissance faciale basée sur le Deep
Learning.
\textbf{Rôle dans le projet} : \\
Elle permet :
\begin{itemize}
  \item l’identification des visages ;
  \item la comparaison faciale ;
  \item la gestion des embeddings faciaux.
\end{itemize}
\textbf{Avantages} : \\
\begin{itemize}
  \item Haute précision ;
  \item Compatible avec les systèmes temps réel ;
  \item Excellentes performances.
\end{itemize}
\begin{figure}[H]
  \centering
  % \includegraphics[width=0.8\textwidth]{images/...}
  \caption{Logo InsightFace.}
\end{figure}
\subsection{Streaming vidéo temps réel}
Définition
Le streaming vidéo temps réel permet la transmission con tinue des flux vidéo avec une faible latence.
\textbf{Rôle dans le projet} : \\
Il permet d’afficher les flux des caméras directement dans la plateforme afin d’assurer une
surveillance instantanée.
\textbf{Avantages} : \\
\begin{itemize}
  \item Surveillance en direct ;
  \item Réactivité rapide ;
  \item Mise à jour instantanée des vidéos.
\end{itemize}
\begin{figure}[H]
  \centering
  % \includegraphics[width=0.8\textwidth]{images/...}
  \caption{Architecture du streaming vidéo temps réel.}
\end{figure}
\subsection{Caméras IP et protocole RTSP}
Définition
Les caméras IP sont des disposi tifs de surveillance connectés au réseau. Le protocole RTSP permet la
transmission des flux vidéo en temps réel.
\textbf{Rôle dans le projet} : \\
Les caméras IP fournissent les flux vidéo u tilisés par les modules d’intelligence ar tificielle pour
l’analyse et la détec tion.
\textbf{Avantages} : \\
\begin{itemize}
  \item Transmission temps réel ;
  \item Bonne qualité vidéo ;
  \item Compatibilité avec plusieurs systèmes.
\end{itemize}
\begin{figure}[H]
  \centering
  % \includegraphics[width=0.8\textwidth]{images/...}
  \caption{Communica tion RTSP entre caméra et serveur.}
\end{figure}
\section{Architecture logicielle}
\subsection{Architecture Client/Serveur}
Définition
L’architecture client/serveur est un modèle informa tique séparant les traitements entre le client et le
serveur.
\textbf{Rôle dans le projet} : \\
Dans notre plateforme :
\begin{itemize}
  \item Angular représente la par tie client ;
  \item Spring Boot représente le serveur principal ;
  \item Python assure les traitements d’intelligence ar tificielle.
\end{itemize}
\textbf{Avantages} : \\
\begin{itemize}
  \item Bonne organisa tion ;
  \item Séparation des responsabilités ;
  \item Facilité de maintenance.
\end{itemize}
\begin{figure}[H]
  \centering
  % \includegraphics[width=0.8\textwidth]{images/...}
  \caption{Architecture client/serveur du système.}
\end{figure}
\subsection{Architecture MVC}
Définition
Le modèle MVC (Model View Controller) est une architecture logicielle perme ttant de séparer les
différentes couches d’une applica tion afin d’améliorer l’organisa tion et la maintenabilité du code.
\textbf{Rôle dans le projet} : \\
Dans notre projet, l’architecture MVC a été adoptée au niveau du Back-End Spring Boot afin de
séparer :
\begin{itemize}
  \item la couche présenta tion ;
  \item la logique mé tier ;
  \item l’accès aux données.
\end{itemize}
\textbf{Structure u tilisée} : \\
\begin{itemize}
  \item Controller : ges tion des requêtes HTTP ;
  \item Service : traitement de la logique mé tier ;
  \item Repository : communica tion avec la base de données ;
  \item Entity : représenta tion des données.
\end{itemize}
\textbf{Avantages} : \\
\begin{itemize}
  \item Bonne organisa tion du code ;
  \item Facilité de maintenance ;
  \item Réutilisation des composants.
\end{itemize}
\begin{figure}[H]
  \centering
  % \includegraphics[width=0.8\textwidth]{images/...}
  \caption{Architecture MVC sous Spring Boot.}
\end{figure}
\subsection{API REST}
Définition
Une API REST est une interface perme ttant la communica tion entre différentes applica tions via le
protocole HTTP.
\textbf{Rôle dans le projet} : \\
Les API REST assurent la communica tion entre Angular, Spring Boot et les différents services de la
plateforme.
\textbf{Avantages} : \\
\begin{itemize}
  \item Simplicité ;
  \item Flexibilité ;
  \item Facilité d’intégra tion.
\end{itemize}
\begin{figure}[H]
  \centering
  % \includegraphics[width=0.8\textwidth]{images/...}
  \caption{Architecture de communica tion REST.}
\end{figure}
\subsection{WebSocket}
Définition
WebSocket est un protocole de communica tion perme ttant des échanges bidirec tionnels en temps
réel entre le client et le serveur.
\textbf{Rôle dans le projet} : \\
WebSocket permet :
\begin{itemize}
  \item l’envoi des alertes en temps réel ;
  \item la mise à jour instantanée des no tifications ;
  \item la synchronisa tion des données de surveillance.
\end{itemize}
\textbf{Avantages} : \\
\begin{itemize}
  \item Communica tion rapide ;
  \item Faible latence ;
  \item Actualisation en temps réel.
\end{itemize}
\begin{figure}[H]
  \centering
  % \includegraphics[width=0.8\textwidth]{images/...}
  \caption{Logo WebSocket.}
\end{figure}
\section{Choix des technologies}
\subsection{Angular}
Définition
Angular est un framework Front- End développé par Google perme ttant de créer des applica tions web
modernes et dynamiques.
\textbf{Rôle dans le projet} : \\
Angular est u tilisé pour développer :
\begin{itemize}
  \item l’interface u tilisateur ;
  \item les tableaux de bord ;
  \item les pages de ges tion ;
  \item les statistiques et visualisa tions.
\end{itemize}
\textbf{Avantages} : \\
\begin{itemize}
  \item Interface moderne ;
  \item Architecture modulaire ;
  \item Bonne performance.
\end{itemize}
\begin{figure}[H]
  \centering
  % \includegraphics[width=0.8\textwidth]{images/...}
  \caption{Logo Angular.}
\end{figure}
\subsection{Spring Boot}
Définition
Spring Boot est un framework Java u tilisé pour le développement des applica tions Back-End et des
API REST.
\textbf{Rôle dans le projet} : \\
Spring Boot assure :
\begin{itemize}
  \item la gestion des utilisateurs ;
  \item les API REST ;
  \item la sécurité ;
  \item la communica tion avec la base de données.
\end{itemize}
\textbf{Avantages} : \\
\begin{itemize}
  \item Développement rapide ;
  \item Sécurité intégrée ;
  \item Compatible avec MySQL.
\end{itemize}
\begin{figure}[H]
  \centering
  % \includegraphics[width=0.8\textwidth]{images/...}
  \caption{Logo Spring Boot.}
\end{figure}
\subsection{Python}
Définition
Python est un langage de programma tion largement u tilisé dans les domaines de l’intelligence
artificielle et du traitement des données.
\textbf{Rôle dans le projet} : \\
Python est u tilisé pour :
\begin{itemize}
  \item l’analyse vidéo ;
  \item la reconnaissance faciale ;
  \item l’intégration des modèles d’intelligence ar tificielle ;
  \item le traitement des flux vidéo.
\end{itemize}
\textbf{Avantages} : \\
\begin{itemize}
  \item Simplicité ;
  \item Large bibliothèque IA ;
  \item Compatible avec OpenCV et YOLOv8.
\end{itemize}
\begin{figure}[H]
  \centering
  % \includegraphics[width=0.8\textwidth]{images/...}
  \caption{Logo Python.}
\end{figure}
\subsection{MySQL}
Définition
MySQL est un système de ges tion de bases de données rela tionnelles open source u tilisé dans les
applications web.
\textbf{Rôle dans le projet} : \\
MySQL permet de stocker :
\begin{itemize}
  \item les employés ;
  \item les départements ;
  \item les rôles ;
  \item les alertes ;
  \item les historiques ;
  \item les informa tions des caméras.
\end{itemize}
\textbf{Avantages} : \\
\begin{itemize}
  \item Bonne performance ;
  \item Facilité d’u tilisation ;
  \item Intégration simple avec Spring Boot.
\end{itemize}
\begin{figure}[H]
  \centering
  % \includegraphics[width=0.8\textwidth]{images/...}
  \caption{Logo MySQL.}
\end{figure}
\subsection{GitHub}
Définition
GitHub est une plateforme de ges tion et d’hébergement de projets u tilisant le système de contrôle de
version Git.
\textbf{Rôle dans le projet} : \\
GitHub a été u tilisé pour :
\begin{itemize}
  \item sauvegarder le code source ;
  \item gérer les versions du projet ;
  \item suivre les modifica tions ;
  \item organiser le développement de la plateforme.
\end{itemize}
\textbf{Avantages} : \\
\begin{itemize}
  \item Gestion des versions ;
  \item Sauvegarde sécurisée ;
  \item Collabora tion simplifiée.
\end{itemize}
\begin{figure}[H]
  \centering
  % \includegraphics[width=0.8\textwidth]{images/...}
  \caption{Logo GitHub.}
\end{figure}
\subsection{Visual Studio Code}
Définition
Visual Studio Code est un éditeur de code source u tilisé pour le développement des applica tions web
et logicielles.
\textbf{Rôle dans le projet} : \\
Il a été utilisé pour développer les différentes par ties de la plateforme.
\textbf{Avantages} : \\
\begin{itemize}
  \item Léger et rapide ;
  \item Compatible avec plusieurs langages ;
  \item Large choix d’extensions.
\end{itemize}
\begin{figure}[H]
  \centering
  % \includegraphics[width=0.8\textwidth]{images/...}
  \caption{Logo Visual Studio Code.}
\end{figure}
\subsection{Thunder Client}
Définition
Thunder Client est une extension légère intégrée à Visual Studio Code perme ttant de tester et de
gérer les API REST directement depuis l’environnement de développement.
\textbf{Rôle dans le projet} : \\
Dans notre projet, Thunder Client a été u tilisé pour tester les différentes API développées avec Spring
Boot, vérifier les requêtes HTTP ainsi que valider les réponses du serveur.
\textbf{Avantages} : \\
\begin{itemize}
  \item Interface simple et intui tive ;
  \item Intégration directe dans Visual Studio Code ;
  \item Facilité de test des API REST ;
  \item Gain de temps durant le développement.
\end{itemize}
\begin{figure}[H]
  \centering
  % \includegraphics[width=0.8\textwidth]{images/...}
  \caption{Logo Thunder Client.}
\end{figure}
\section{Conclusion}
Dans ce chapitre, nous avons présenté le domaine d’applica tion de notre projet ainsi que les
principaux concepts, architectures, technologies et ou tils utilisés dans le développement de notre
plateforme intelligente de surveillance et de sécurité.
Les technologies choisies, telles qu’Angular, Spring Boot, Python, MySQL, YOLOv8, OpenCV et
InsightFace, perme ttent de développer une solu tion performante, intelligente et adaptée aux besoins
des entreprises modernes.
Dans le chapitre suivant, nous présenterons l’analyse et la spécifica tion des besoins fonc tionnels et
non fonctionnels du système ainsi que les différents diagrammes de concep tion.



% ===========================================================
% CHAPITRE 3 : SPÉCIFICATION ET CONCEPTION
% ===========================================================
\chapter{Spécification des besoins et Conception}

\section{Introduction}
Ce chapitre détaille les fonctionnalités attendues du système Farm AI ainsi que la modélisation
UML globale et les diagrammes d'architecture.

\section{Spécification des besoins}

\subsection{Besoins fonctionnels}
Le système doit répondre aux fonctionnalités suivantes :

\subsection{Besoins fonctionnels et non fonctionnels}

Le tableau~\ref{tab:besoins_fonctionnels} récapitule l'ensemble des besoins fonctionnels
identifiés, classés par module et par acteur concerné.

\begin{longtable}{|c|>{\raggedright\arraybackslash}p{5cm}|>{\raggedright\arraybackslash}p{6cm}|l|}
  \caption{Besoins fonctionnels du système Farm AI} \label{tab:besoins_fonctionnels} \\
  \hline
  \rowcolor{gray!20}
  \textbf{ID} & \textbf{Besoin fonctionnel} & \textbf{Description} & \textbf{Acteur} \\
  \hline
  \endfirsthead
  \hline
  \rowcolor{gray!20}
  \textbf{ID} & \textbf{Besoin fonctionnel} & \textbf{Description} & \textbf{Acteur} \\
  \hline
  \endhead
  \hline
  \endfoot

  BF01 & Authentification classique &
    Se connecter via email et mot de passe sécurisé avec génération de token JWT. &
    Tous \\
  \hline
  BF02 & Authentification biométrique &
    Se connecter par reconnaissance faciale en temps réel sans saisie de mot de passe. &
    Tous \\
  \hline
  BF03 & Gérer les employés &
    Ajouter, modifier, supprimer et consulter les fiches des employés avec leur métier, email et statut. &
    Admin \\
  \hline
  BF04 & Enregistrer un Face ID employé &
    Capturer et enregistrer l'encodage biométrique du visage d'un employé pour le pointage automatique. &
    Admin \\
  \hline
  BF05 & Approuver/Rejeter un employé &
    Valider ou refuser un employé nouvellement créé avant son affectation à un département. &
    Admin \\
  \hline
  BF06 & Gérer les gestionnaires &
    Créer, modifier et supprimer des comptes gestionnaires. Un email d'activation avec mot de passe est envoyé automatiquement via Gmail. &
    Admin \\
  \hline
  BF07 & Gérer les observateurs &
    Créer, modifier et supprimer des comptes observateurs. Un email d'activation est envoyé automatiquement via Gmail. &
    Admin \\
  \hline
  BF08 & Gérer les départements &
    Créer, modifier et supprimer des départements avec nom, horaires, responsable (manager) et besoins en personnel par métier. &
    Admin \\
  \hline
  BF09 & Activer un compte par email &
    Cliquer sur le lien d'activation reçu par Gmail pour activer son compte et enregistrer son visage. &
    Manager / Viewer \\
  \hline
  BF10 & Enregistrer son visage (onboarding) &
    Enregistrer ses données biométriques lors de la première connexion après activation du compte. &
    Manager / Viewer \\
  \hline
  BF11 & Gérer les caméras &
    Consulter la liste des caméras, vérifier leur statut (active/inactive) et déclencher l'analyse IA. &
    Admin \\
  \hline
  BF12 & Lancer la surveillance IA &
    Activer l'analyse intelligente sur une caméra : détection de personnes (YOLO), reconnaissance faciale (FaceNet) et analyse vestimentaire. &
    Admin \\
  \hline
  BF13 & Pointage automatique &
    Enregistrer automatiquement les entrées et sorties des employés identifiés avec horodatage précis. &
    Service IA \\
  \hline
  BF14 & Consulter le tableau de bord &
    Visualiser les KPIs (présents, heures travaillées, entrées, caméras actives), les graphiques de présence et les insights IA. &
    Admin \\
  \hline
  BF15 & Gérer les alertes &
    Consulter les alertes générées lors de détections d'anomalies (personne non identifiée, intrusion). &
    Admin \\
  \hline
  BF16 & Suggestion d'équipe &
    Recevoir une proposition automatique d'équipe basée sur les besoins en personnel du département et la disponibilité des employés. &
    Manager \\
  \hline
  BF17 & Valider/Modifier l'équipe &
    Accepter la suggestion d'équipe ou la modifier manuellement avant validation. Un email de confirmation est envoyé à l'admin et aux employés. &
    Manager \\
  \hline
  BF18 & Tableau de bord manager &
    Consulter un tableau de bord dédié avec les employés du département, les caméras associées et les alertes. &
    Manager \\
  \hline
  BF19 & Tableau de bord observateur &
    Consulter en lecture seule la liste des employés et des départements. &
    Viewer \\
  \hline
  BF20 & Notifications par email (Gmail) &
    Envoyer automatiquement des emails via Gmail SMTP pour : création de compte, activation, affectation, validation d'équipe et suppression. &
    Système \\
  \hline
  BF21 & Gérer son compte &
    Modifier ses informations personnelles, changer son mot de passe et gérer ses données biométriques. &
    Tous \\
  \hline
  BF22 & Changer la langue &
    Basculer l'interface entre le français et l'anglais. &
    Tous \\
  \hline
\end{longtable}

Le tableau~\ref{tab:besoins_non_fonctionnels} présente les besoins non fonctionnels qui ont
guidé les choix techniques tout au long du développement.

\begin{table}[H]
  \centering
  \caption{Besoins non fonctionnels du système Farm AI}
  \label{tab:besoins_non_fonctionnels}
  \renewcommand{\arraystretch}{1.4}
  \begin{tabularx}{\textwidth}{|c|l|X|}
    \hline
    \rowcolor{gray!20}
    \textbf{ID} & \textbf{Catégorie} & \textbf{Description} \\
    \hline
    BNF01 & Performance &
      Le traitement du flux vidéo doit s'effectuer avec une latence inférieure à 2 secondes entre la capture et l'affichage du résultat. \\
    \hline
    BNF02 & Sécurité &
      Toutes les requêtes sont authentifiées par JWT. Les données biométriques sont stockées sous forme d'encodages mathématiques. Les mots de passe sont hachés avec BCrypt. \\
    \hline
    BNF03 & Fiabilité IA &
      Le taux d'identification correcte doit être supérieur à 90\,\%. Le modèle FaceNet retenu atteint 97,2\,\%. \\
    \hline
    BNF04 & Ergonomie &
      L'interface doit être intuitive pour un utilisateur non technicien. Toutes les actions critiques sont confirmées par une fenêtre de validation. \\
    \hline
    BNF05 & Évolutivité &
      L'architecture micro-services permet de faire évoluer chaque composant indépendamment. \\
    \hline
    BNF06 & Contrôle d'accès &
      Le système applique un contrôle d'accès basé sur les rôles (RBAC) avec redirection automatique selon le rôle de l'utilisateur connecté. \\
    \hline
    BNF07 & Disponibilité &
      Le service de notification par email doit fonctionner de manière asynchrone pour ne pas bloquer les opérations principales. \\
    \hline
    BNF08 & Internationalisation &
      L'interface supporte le français et l'anglais avec basculement dynamique sans rechargement. \\
    \hline
    BNF09 & Compatibilité &
      L'application est compatible avec les navigateurs modernes (Chrome, Firefox, Edge). \\
    \hline
  \end{tabularx}
\end{table}



\section{Conception globale et modélisation UML}



\subsection{Gérer les gestionnaires}

\begin{table}[H]
  \centering
  \caption{Description textuelle du cas d'utilisation « Gérer les gestionnaires »}
  \label{tab:cu_managers}
  \renewcommand{\arraystretch}{1.4}
  \begin{tabularx}{\textwidth}{|>{\bfseries}l|L|}
    \hline
    \rowcolor{gray!20}
    \multicolumn{2}{|c|}{\textbf{Créer un gestionnaire}} \\
    \hline
    Nom C.U & Créer un gestionnaire \\
    \hline
    Acteur & Administrateur \\
    \hline
    Description & L'administrateur crée un compte gestionnaire. Un mot de passe est généré
      automatiquement et un email d'activation est envoyé via Gmail. \\
    \hline
    Précondition & L'administrateur est authentifié et se trouve sur le module de gestion
      des gestionnaires. \\
    \hline
    Postcondition & Le compte gestionnaire est créé en base de données avec le rôle
      ROLE\_MANAGER, le statut désactivé, et un email d'activation est envoyé. \\
    \hline
    Scénario nominal &
      1. L'administrateur saisit le prénom, nom, email et téléphone.\newline
      2. Il clique sur « Enregistrer ».\newline
      3. Le backend génère un mot de passe aléatoire et un token d'activation.\newline
      4. Le compte est créé avec \texttt{enabled = false}.\newline
      5. Un email contenant le mot de passe et le lien d'activation est envoyé via Gmail SMTP.\newline
      6. Le gestionnaire apparaît dans la liste avec le statut « En attente d'activation ». \\
    \hline
    Scénarios alternatifs &
      -- Si l'email existe déjà, une erreur 400 est retournée.\newline
      -- Si le service email est indisponible, le compte est créé mais l'email n'est pas envoyé. \\
    \hline
  \end{tabularx}
\end{table}

\subsection{Gérer les départements}

\begin{table}[H]
  \centering
  \caption{Description textuelle du cas d'utilisation « Créer un département »}
  \label{tab:cu_departments}
  \renewcommand{\arraystretch}{1.4}
  \begin{tabularx}{\textwidth}{|>{\bfseries}l|L|}
    \hline
    \rowcolor{gray!20}
    \multicolumn{2}{|c|}{\textbf{Créer un département}} \\
    \hline
    Nom C.U & Créer un département \\
    \hline
    Acteur & Administrateur \\
    \hline
    Description & L'administrateur crée un département en définissant son nom, ses horaires
      de travail, le gestionnaire responsable et les besoins en personnel par métier. \\
    \hline
    Précondition & L'administrateur est authentifié. Au moins un gestionnaire existe
      dans le système. \\
    \hline
    Postcondition & Le département est enregistré en base avec les besoins en docteurs,
      électriciens et ouvriers. \\
    \hline
    Scénario nominal &
      1. L'administrateur saisit le nom du département.\newline
      2. Il sélectionne un gestionnaire responsable dans la liste déroulante.\newline
      3. Il définit les horaires d'ouverture et de fermeture.\newline
      4. Il précise les besoins en personnel : nombre de docteurs, d'électriciens et d'ouvriers requis.\newline
      5. Il clique sur « Enregistrer ».\newline
      6. Le département est créé et affiché sous forme de carte avec les détails. \\
    \hline
    Scénarios alternatifs &
      -- Si aucun gestionnaire n'est sélectionné, un avertissement s'affiche.\newline
      -- Si les nombres d'employés sont négatifs, une validation empêche la soumission. \\
    \hline
  \end{tabularx}
\end{table}

\subsection{Activer un compte et enregistrer son visage}

\begin{table}[H]
  \centering
  \caption{Description textuelle du cas d'utilisation « Activer un compte »}
  \label{tab:cu_activate}
  \renewcommand{\arraystretch}{1.4}
  \begin{tabularx}{\textwidth}{|>{\bfseries}l|L|}
    \hline
    \rowcolor{gray!20}
    \multicolumn{2}{|c|}{\textbf{Activer un compte par email}} \\
    \hline
    Nom C.U & Activer un compte par email \\
    \hline
    Acteur & Gestionnaire / Observateur \\
    \hline
    Description & Le gestionnaire ou l'observateur clique sur le lien d'activation reçu
      par email pour activer son compte, puis enregistre son visage. \\
    \hline
    Précondition & Un email d'activation a été envoyé. Le token est valide et non expiré. \\
    \hline
    Postcondition & Le compte est activé (\texttt{enabled = true}) et l'utilisateur est
      redirigé vers la page d'enregistrement facial. \\
    \hline
    Scénario nominal &
      1. L'utilisateur clique sur le lien d'activation dans l'email reçu.\newline
      2. L'application Angular extrait le token de l'URL.\newline
      3. Une requête POST est envoyée au backend pour valider le token.\newline
      4. Le backend active le compte et retourne un token JWT.\newline
      5. L'utilisateur est redirigé vers la page « Face Setup ».\newline
      6. Il enregistre son visage via la caméra.\newline
      7. Après succès, il est redirigé vers son tableau de bord dédié. \\
    \hline
    Scénarios alternatifs &
      -- Si le token est invalide ou expiré, un message d'erreur s'affiche.\newline
      -- Si la caméra est indisponible, l'enregistrement facial échoue. \\
    \hline
  \end{tabularx}
\end{table}

\subsection{Suggestion et validation d'équipe}

\begin{table}[H]
  \centering
  \caption{Description textuelle du cas d'utilisation « Valider une équipe »}
  \label{tab:cu_team}
  \renewcommand{\arraystretch}{1.4}
  \begin{tabularx}{\textwidth}{|>{\bfseries}l|L|}
    \hline
    \rowcolor{gray!20}
    \multicolumn{2}{|c|}{\textbf{Valider une équipe}} \\
    \hline
    Nom C.U & Valider une équipe \\
    \hline
    Acteur & Gestionnaire \\
    \hline
    Description & Le gestionnaire consulte la proposition automatique d'équipe générée
      par le système, la modifie si nécessaire, puis la valide pour affecter les employés
      à son département. \\
    \hline
    Précondition & Le gestionnaire est authentifié. Son département a des besoins en
      personnel définis. Des employés disponibles et approuvés existent en base. \\
    \hline
    Postcondition & Les employés sélectionnés sont affectés au département. Des emails
      de confirmation sont envoyés au gestionnaire, à l'administrateur et aux employés. \\
    \hline
    Scénario nominal &
      1. Le gestionnaire accède à son tableau de bord.\newline
      2. Le système affiche une suggestion d'équipe basée sur les besoins du département
         et la disponibilité des employés (visage enregistré, statut APPROVED, non assigné).\newline
      3. Le gestionnaire peut ajouter ou retirer des employés de la sélection.\newline
      4. Il clique sur « Valider l'équipe ».\newline
      5. Le backend vérifie que chaque employé est disponible et validé.\newline
      6. Les employés sont assignés au département (\texttt{available = false}).\newline
      7. Un email de confirmation est envoyé au gestionnaire, à l'admin et à chaque employé. \\
    \hline
    Scénarios alternatifs &
      -- Si un employé n'est pas approuvé ou n'a pas de visage enregistré, la validation échoue.\newline
      -- Le gestionnaire peut modifier l'équipe après validation initiale. \\
    \hline
  \end{tabularx}
\end{table}

\subsection{Architecture globale du système}
Le système adopte une architecture micro-services séparant :
\begin{itemize}
  \item \textbf{Service IA} (Python/FastAPI) : traitement vidéo et IA ;
  \item \textbf{Backend} (Java/Spring Boot) : logique métier et persistance ;
  \item \textbf{Frontend} (Angular/TypeScript) : interface utilisateur.
\end{itemize}

\section{Conclusion}
Ce chapitre a établi l'ensemble des besoins du système et défini son architecture. Le chapitre
suivant présentera les choix technologiques retenus.

% ===========================================================
% CHAPITRE 4 : ENVIRONNEMENT ET CHOIX TECHNIQUES
% ===========================================================
\chapter{Environnement et Choix Techniques}

\section{Introduction}
Ce chapitre expose l'environnement matériel et les différentes technologies employées dans la
réalisation de la plateforme Farm AI.

\section{Architecture logicielle}
L'architecture multi-tiers adoptée connecte le flux IA (Python/FastAPI), le serveur central
(Java/Spring Boot) et l'interface web (Angular) via des API RESTful sécurisées par JWT.

\section{Environnement matériel et outils}

\subsection{Environnement matériel}
Le développement a été effectué sur des PC performants équipés de :
\begin{itemize}
  \item GPU NVIDIA pour l'accélération du traitement vidéo ;
  \item RAM $\geq$ 16 GB pour le multitâche ;
  \item SSD pour les performances d'E/S.
\end{itemize}

\subsection{Outils de développement}
L'utilisation de VS Code, IntelliJ IDEA et Git a permis une structure efficace du travail et du
contrôle de versions.

\section{Technologies de l'IA (Service Python)}

\subsection{Python et FastAPI}
Python a été privilégié pour sa variété de bibliothèques IA, FastAPI pour l'exposition ultra-rapide
des résultats sous forme d'API REST.

\begin{infobox}
  FastAPI offre une performance comparable à Go et NodeJS avec la facilité de Python.
\end{infobox}

\subsection{OpenCV et modèles de Computer Vision}
\begin{itemize}
  \item \textbf{OpenCV} : capture et manipulation des frames vidéo ;
  \item \textbf{dlib} : encodage précis des visages (Face recognition) ;
  \item \textbf{YOLO v8} : détection de personnes et d'objets en temps réel.
\end{itemize}

\section{Technologies Backend (Java)}

\subsection{Service de notification par email (Gmail SMTP)}

Le système Farm AI intègre un service de notification par email utilisant le protocole
\textbf{Gmail SMTP}. Ce service, implémenté dans le backend Spring Boot via la classe
\texttt{EmailService}, utilise l'API \texttt{JavaMailSender} de Spring pour envoyer des
emails de manière asynchrone (annotation \texttt{@Async}).

La configuration dans le fichier \texttt{application.properties} utilise le serveur
SMTP de Google (\texttt{smtp.gmail.com}, port 587) avec l'authentification TLS activée.
Un mot de passe d'application Google est utilisé pour sécuriser la connexion.

Le service gère huit types de notifications automatiques :

\begin{enumerate}
  \item \textbf{Création d'un compte employé :} notification informant l'employé que son
    profil a été créé dans le système.
  \item \textbf{Validation d'un employé :} email de bienvenue envoyé après approbation par
    l'administrateur.
  \item \textbf{Activation de compte gestionnaire :} email contenant le mot de passe
    généré et un lien d'activation unique avec token.
  \item \textbf{Activation de compte observateur :} même mécanisme que pour les gestionnaires,
    avec un lien d'activation dédié.
  \item \textbf{Affectation à un département :} notification envoyée à l'employé lors de
    son assignation, incluant le nom du département, le responsable et les horaires.
  \item \textbf{Validation d'équipe (vers l'admin) :} récapitulatif envoyé à l'administrateur
    lorsqu'un gestionnaire valide la composition de son équipe.
  \item \textbf{Confirmation d'équipe (vers le manager) :} confirmation envoyée au
    gestionnaire après validation réussie.
  \item \textbf{Suppression de département :} notification informant le gestionnaire
    de la suppression de son département.
\end{enumerate}

\subsection{Java et Spring Boot}
Le socle applicatif principal garantit la sécurité avec Spring Security et JWT. Spring Boot
facilite le déploiement et la gestion des dépendances.

\subsection{MySQL (correction)}

MySQL est la base de données relationnelle utilisée pour la persistance de toutes
les données du système : utilisateurs (administrateurs, gestionnaires, observateurs),
employés, départements, pointages, caméras et alertes. La connexion est configurée
dans le fichier \texttt{application.properties} avec le driver JDBC MySQL sur le
port 3306, et Hibernate gère automatiquement le schéma via \texttt{ddl-auto=update}.
Une base de données relationnelle robuste pour la sauvegarde de l'historique des présences,
des alertes et des statistiques.

\section{Technologies Frontend (Angular)}

\subsection{Framework Angular et TypeScript}
Angular garantit une application mono-page (SPA) réactive avec un typage fort via TypeScript.

\subsection{Design et SCSS}
Le design adopte des thèmes sombres et une approche glassmorphism pour le rendu cyber-technologique.

\section{Conclusion}
Ce chapitre a justifié les choix techniques. Le chapitre suivant présentera la réalisation concrète
et les résultats des tests.

% ===========================================================
% CHAPITRE 5 : RÉALISATION ET TESTS
% ===========================================================
\chapter{Réalisation et Tests}

\section{Introduction}
Cette dernière étape permet de visualiser les interfaces développées et d'analyser les résultats
des tests de l'IA.

\section{Présentation des interfaces}

\subsection{Interface d'authentification}
L'interface de connexion offre deux modes d'authentification :
\begin{enumerate}
  \item Authentification classique (email/mot de passe) ;
  \item Authentification biométrique (reconnaissance faciale).
\end{enumerate}

\subsection{Tableau de bord analytique}
Le tableau de bord affiche :
\begin{itemize}
  \item Indicateurs clés (KPIs) ;
  \item Graphiques de présence ;
  \item Alertes actives ;
  \item Timeline des activités.
\end{itemize}

\subsection{Gestion des employés}
Module CRUD complet pour :
\begin{itemize}
  \item Ajouter/modifier/supprimer des employés ;
  \item Enregistrer le Face ID ;
  \item Assigner des rôles et permissions.
\end{itemize}

\subsection{Gestion des caméras}
Contrôle de la surveillance intelligente avec :
\begin{itemize}
  \item Liste des caméras connectées ;
  \item Statut de connexion en temps réel ;
  \item Déclenchement de l'analyse IA ;
  \item Visualisation des alertes.
\end{itemize}

\section{Tests et validation}

\subsection{Tests de précision de l'IA}
Les performances des modèles sous diverses conditions d'éclairage :

\begin{table}[H]
  \centering
  \caption{Comparaison des performances des modèles IA}
  \label{tab:perf_ia}
  \begin{tabular}{lccc}
    \toprule
    \textbf{Modèle} & \textbf{Précision (\%)} & \textbf{Rappel (\%)} & \textbf{F1-Score} \\
    \midrule
    dlib ResNet & 94,5 & 92,0 & 0,932 \\
    FaceNet & 97,2 & 95,8 & 0,965 \\
    YOLO v8 & 96,8 & 94,5 & 0,956 \\
    \bottomrule
  \end{tabular}
\end{table}

\subsection{Tests de charge}
Le service FastAPI s'est avéré très résilient face à plusieurs flux vidéo simultanés. Les tests
ont confirmé la capacité du système à gérer 10+ flux en parallèle sans dégradation.

\subsection{Tests de l'interface utilisateur}
Les interactions sur le tableau de bord Angular sont fluides et réagissent bien aux mises à jour
en temps réel via WebSocket.

\section{Conclusion}
Les tests réalisés confirment la robustesse et la précision du système Farm AI. L'ensemble des
objectifs fixés en début de projet ont été atteints.

% ===========================================================
% CONCLUSION GÉNÉRALE ET PERSPECTIVES
% ===========================================================
\chapter*{Conclusion générale et Perspectives}
\addcontentsline{toc}{chapter}{Conclusion générale}

Ce projet a permis de concevoir et de développer de bout en bout la plateforme \textbf{Farm AI},
un système de gestion agricole intelligent intégrant :

\begin{itemize}
  \item La reconnaissance faciale en temps réel ;
  \item Le suivi automatique des présences ;
  \item Un tableau de bord analytique avancé ;
  \item Une architecture scalable et sécurisée.
\end{itemize}

Les objectifs fonctionnels et non fonctionnels définis en amont ont été atteints, et les tests
de validation ont confirmé la fiabilité du système.

\subsection*{Perspectives d'évolution}
\begin{itemize}
  \item Intégration de caméras thermiques pour l'évaluation de la santé ;
  \item Déploiement sur Cloud (AWS/Azure) avec Kubernetes ;
  \item Application mobile pour les gestionnaires nomades ;
  \item Extension du modèle IA pour la détection d'équipements de sécurité ;
  \item Intégration de modèles de prédiction (machine learning avancé).
\end{itemize}

% ===========================================================
% BIBLIOGRAPHIE
% ===========================================================
\begin{thebibliography}{99}
\addcontentsline{toc}{chapter}{Bibliographie}

\bibitem{angular}
Angular Development Team.
\textit{Angular Documentation : The modern web application framework}.
[En ligne]. Disponible sur : \url{https://angular.io/}
(Consulté le 01 juin 2025)

\bibitem{opencv}
Bradski, G. et Kaehler, A.
\textit{Learning OpenCV 3: Computer Vision in C++ with the OpenCV Library}.
O'Reilly Media, 2016, 1000 p.
[En ligne]. Disponible sur : \url{https://opencv.org/}

\bibitem{springboot}
Spring Framework Team.
\textit{Spring Boot Reference Documentation}.
[En ligne]. Disponible sur : \url{https://spring.io/projects/spring-boot}
(Consulté le 01 juin 2025)

\bibitem{yolo}
Redmon, J. et Farhadi, A.
\textit{YOLOv3: An Incremental Improvement}.
arXiv:1804.02767, 2018.

\bibitem{fastapi}
Ramirez, S.
\textit{FastAPI Framework - Build fast APIs}.
[En ligne]. Disponible sur : \url{https://fastapi.tiangolo.com/}
(Consulté le 01 juin 2025)

\bibitem{dlib}
King, D. E.
\textit{Dlib-ml: A Machine Learning Toolkit}.
Journal of Machine Learning Research, 2009.

\end{thebibliography}

% ===========================================================
% ANNEXES
% ===========================================================
\appendix
\chapter{Annexes}

\section*{Annexe A : Listing des scripts principaux}
\addcontentsline{toc}{section}{Annexe A : Listing des scripts}

\subsection*{A.1. Script Python : camera\_ai\_stream.py}
Script de traitement vidéo en temps réel avec intégration YOLO et reconnaissance faciale.

\subsection*{A.2. Configuration Spring Boot}
Fichier application.properties avec la configuration de la base de données et JWT.

\section*{Annexe B : Diagrammes UML complets}
\addcontentsline{toc}{section}{Annexe B : Diagrammes UML}

\subsection*{B.1. Diagramme de classes}
Structure complète du modèle de données partagé entre applications.

\subsection*{B.2. Diagramme des cas d'utilisation}
Interactions complètes entre acteurs et système.

\section*{Annexe C : Schéma de la base de données}
\addcontentsline{toc}{section}{Annexe C : Schéma base de données}

Schéma entité-association (ERD) montrant les tables principales et leurs relations.

\end{document}
